\documentclass[12pt]{article}

\usepackage[margin=1in]{geometry}
\usepackage{amsmath,amssymb}
\usepackage[T1]{fontenc}
\usepackage{lmodern}

\title{Homework 1}
\author{Hanvesh Aduru}
\date{\today}

\begin{document}

\maketitle

\section*{1. Factorization and Divisibility}

\subsection*{(a)}

Factor
\[
2^{15}-1=32,767.
\]

Using the difference of cubes,
\[
a^3-b^3=(a-b)(a^2+ab+b^2),
\]
we get
\[
2^{15}-1=(2^5)^3-1^3.
\]

Therefore,
\[
(2^5-1)(2^{10}+2^5+1).
\]

Since
\[
2^5-1=31
\]
and
\[
2^{10}+2^5+1=1024+32+1=1057,
\]
we have
\[
\boxed{32,767=31\cdot1057}.
\]

\subsection*{(b)}

Since
\[
2^{15}-1=32,767
\]
divides
\[
2^{32,767}-1,
\]
we can choose
\[
\boxed{x=32,767}.
\]

This satisfies
\[
1<x<2^{32,767}-1.
\]

\section*{2. Conjectures}

Consider the following values:

\[
\begin{array}{c|c|c|c|c}
n & 3^n-1 & \text{Prime?} & 3^n-2^n & \text{Prime?}\\
\hline
1 & 2 & \text{Yes} & 1 & \text{No}\\
2 & 8 & \text{No} & 5 & \text{Yes}\\
3 & 26 & \text{No} & 19 & \text{Yes}\\
4 & 80 & \text{No} & 65 & \text{No}\\
5 & 242 & \text{No} & 211 & \text{Yes}\\
6 & 728 & \text{No} & 665 & \text{No}
\end{array}
\]

A conjecture is that for $n>1$,
\[
3^n-1
\]
is never prime. This is because $3^n$ is odd, so $3^n-1$ is even and greater than $2$.

Also, $3^n-2^n$ can be prime for some values of $n$, but not for every value of $n$.

If $n$ is composite, then $3^n-2^n$ is also composite because expressions of the form
\[
a^n-b^n
\]
can be factored when $n$ is composite.

\section*{3. Finding a New Prime}

The method is to multiply the given primes together and add $1$.

\subsection*{(a)}

For the primes $2,3,5,$ and $7$:

\[
2\cdot3\cdot5\cdot7+1=211.
\]

Since $211$ is prime and is not in the original list,

\[
\boxed{211}
\]

is the required prime.

\subsection*{(b)}

For $2,5,$ and $11$:

\[
2\cdot5\cdot11+1=111.
\]

However,
\[
111=3\cdot37,
\]
so $111$ is not prime.

Including $3$ and trying again gives

\[
2\cdot5\cdot11\cdot3+1=331.
\]

Since $331$ is prime,

\[
\boxed{331}
\]

is a prime different from $2,5,$ and $11$.

\section*{4. Five Consecutive Non-Prime Integers}

Five consecutive integers that are not prime are

\[
\boxed{24,25,26,27,28}.
\]

Indeed,

\[
24=2\cdot12,
\]
\[
25=5\cdot5,
\]
\[
26=2\cdot13,
\]
\[
27=3\cdot9,
\]
and
\[
28=2\cdot14.
\]

Therefore none of these five integers is prime.

\section*{5. Perfect Numbers}

The first perfect numbers are

\[
6,\quad28,\quad496,\quad8128.
\]

Therefore, the two additional perfect numbers are

\[
\boxed{496\text{ and }8128}.
\]

For $496$,

\[
1+2+4+8+16+31+62+124+248=496.
\]

For $8128$,

\[
1+2+4+8+16+32+64+127+254
+508+1016+2032+4064=8128.
\]

\section*{6. Triplet Primes}

The prime triplets in which each pair of adjacent numbers differs by $2$ are

\[
\boxed{(3,5,7)}
\]

and

\[
\boxed{(5,7,11)}.
\]

These are the only two such triplets.

\section*{7. Amicable Numbers}

We show that $(220,284)$ is an amicable pair.

The positive divisors of $220$ that are smaller than $220$ are

\[
1,2,4,5,10,11,20,22,44,55,110.
\]

Their sum is

\[
1+2+4+5+10+11+20+22+44+55+110
=284.
\]

Now consider $284$. Its positive divisors smaller than $284$ are

\[
1,2,4,71,142.
\]

Their sum is

\[
1+2+4+71+142=220.
\]

Thus, the sum of the proper divisors of $220$ is $284$, and the sum of the proper divisors of $284$ is $220$.

Therefore,

\[
\boxed{(220,284)\text{ is an amicable pair}.}
\]

\end{document}